\documentclass[11pt]{article}
\usepackage[T1]{fontenc}
\usepackage{textcomp}
\usepackage[margin=0.75in]{geometry}
\usepackage{amsmath,amssymb,amsthm}
\usepackage{array,booktabs}
\usepackage{hyperref}
\setlength{\parindent}{0pt}
\newtheorem{theorem}{Theorem}
\theoremstyle{definition}
\newtheorem{definition}{Definition}
\title{Flow Proof Artifact: Nat-mul-one-three}
\author{Generated by \texttt{flow doc proof}}
\date{Flow Proof Book}
\begin{document}
\maketitle
One times three is three for natural numbers.\\[0.5em]
\textbf{Source.} peano — https://en.wikipedia.org/wiki/Peano\_axioms\\[0.5em]
\bigskip
\noindent{\large \textbf{Derived fact 1.} \textit{1 * 3 = 3} \hfill the natural numbers \cdot multiplication \cdot one times three is three}\par
\smallskip\noindent\textit{Coordinate: the natural numbers \cdot multiplication \cdot one times three is three}\par
\smallskip\noindent\textit{Source: peano}\par
\smallskip\noindent\textit{Built on: one is the left identity, for multiplication on the natural numbers}\par
\medskip
\noindent\textbf{Goal.} 1 * 3 = 3\par
\noindent$\displaystyle \mathrm{succ}(0) * three = three$\par
\medskip
\noindent
\begin{tabular}{@{} >{\raggedright\arraybackslash}p{0.06\textwidth} >{\raggedright\arraybackslash}p{0.47\textwidth} >{\raggedleft\arraybackslash}p{0.41\textwidth} @{}}
\textbf{\#} & \textbf{Proof} & \textbf{Mathematics} \\
\midrule
\textbf{1.} & We prove that one times three is three for multiplication on the natural numbers. & \\[0.45em]
\textbf{2.} & Let three = succ(succ(succ(0))). & \\[0.45em]
\textbf{3.} & We invoke the derived fact governing multiplication on the natural numbers: one is the left identity, for multiplication on the natural numbers (instantiated for three). & \\[0.45em]
\textbf{4.} & From step 2 and step 3, this implies the successor of 0 times three equals three. Hence proven. & $\displaystyle \mathrm{succ}(0) * three = three$ \\[0.45em]
\bottomrule
\end{tabular}
\label{thm:Nat-multiplication-one_times_three_is_three}
\par\medskip\hrule
\end{document}
